\chapter{Literature Review}\label{ch:literature}

This chapter places the project in the existing body of knowledge. It moves from the business case for
analytics, through the process models and techniques used to build analytical systems, to the specific
methods applied in this work. It closes by identifying the gap that the project addresses.

\section{Business intelligence and analytics}
\citet{chen2012} describe the evolution of business intelligence and analytics (BI\&A) in three stages. The
first, BI\&A~1.0, centred on data warehouses and structured reporting. The second added web and unstructured
data, and the third added mobile and sensor data. They argue that the main value lies in turning data into
\emph{decisions}, not merely into reports. \citet{davenport2007} reached a similar conclusion from
organisational case studies: firms that ``compete on analytics'' use statistical and predictive models
routinely in their core processes rather than occasionally in projects. \citet{sharda2018} organise analytics
into three levels. \emph{Descriptive} analytics asks what happened, \emph{predictive} analytics asks what will
happen, and \emph{prescriptive} analytics asks what should be done. Most organisations remain at the descriptive
level. The three management questions in the project proposal map directly onto this hierarchy, and the
framework built here covers the descriptive and predictive levels while giving rule-based prescriptions.

\citet{provost2013} emphasise \emph{data-analytic thinking}. Analytics is valuable only when a business problem
has been framed precisely enough to be answered by data, and models must be judged by their effect on
decisions, not by technical accuracy alone. \citet{shmueli2010} adds that explanatory and predictive modelling
are different activities with different goals and validation methods. A model that explains past variation
well may forecast poorly, and the reverse can also be true. This distinction shapes the present study. The
explanatory analyses (EDA, hypothesis tests, elasticity regression) are kept separate from the predictive
analyses (out-of-sample forecasting and classification), and each is evaluated on its own terms.

\section{Process models for analytics projects}
The Cross-Industry Standard Process for Data Mining (CRISP-DM) remains the most widely used methodology for
analytics projects \citep{shearer2000}. Its six phases are business understanding, data understanding, data
preparation, modelling, evaluation and deployment. They are iterative, and the model explicitly recognises that
data preparation usually takes most of the effort. The ten-phase plan in the project proposal is a refinement
of CRISP-DM. It adds explicit phases for feature engineering, dashboard development, system evaluation and
documentation. \citet{han2012} describe a comparable knowledge-discovery process (cleaning, integration,
selection, transformation, mining, evaluation and presentation), and the framework in \cref{ch:system} follows
both.

\section{Data quality and preprocessing}
Real-world data are rarely ready for analysis. \citet{garcia2015} classify preprocessing into data cleaning
(missing values, noise and outliers), integration, transformation (normalisation and discretisation) and
reduction. They stress that the choice of technique should depend on the mechanism that produced the problem.
For missing values, deletion is safe only when values are missing completely at random and few in number.
Otherwise, imputation is preferable, and imputation from deterministic relationships between columns is the
most accurate option of all. Textbooks on data mining make the same point. \citet{tan2019} and
\citet{witten2017} devote substantial space to data quality before any modelling, and \citet{kuhn2019} show
that the way predictors are constructed and encoded often matters more for accuracy than the choice of model.
\citet{wickham2014} adds the principle of \emph{tidy data}: each variable forms a
column, each observation a row, and each type of observational unit a table. This principle underlies the
star-schema design used here.

Outlier treatment requires particular care in business data. Statistical rules such as the $z$-score assume
approximately normal data and are distorted by the very outliers they look for. The inter-quartile range (IQR)
rule and the median absolute deviation (MAD) are more robust \citep{han2012}. \citet{aggarwal2017} reviews
outlier analysis in depth and stresses that the definition of an outlier depends on the application and on the
model of ``normal'' behaviour that is assumed. Not every outlier is an error,
however. A bulk order from a distributor is unusual but genuine, and deleting it would bias demand estimates
downwards. The approach taken here therefore distinguishes \emph{impossible} values (domain-rule violations,
which are corrected) from \emph{unusual} values (statistical outliers, which are flagged and kept).

A further, often under-estimated risk is \emph{data leakage}: information that would not be available at
prediction time slips into the training data. \citet{kaufman2012} show that leakage is common in published
studies and competitions and that it produces models which look excellent in testing but fail in deployment.
Typical sources are features computed from the target itself, rolling statistics that include the current
period, and preprocessing parameters (for example, scaling constants or target encodings) estimated on the
whole dataset instead of the training portion only. All three are avoided by design in this project.

\section{Data warehousing}
\citet{kimball2013} established dimensional modelling as the standard design for analytical databases. A
central \emph{fact table} stores measurable events (here, sales transactions). Surrounding \emph{dimension
tables} describe the context of each event (date, product, customer segment). The resulting star schema is easy
to understand, efficient for aggregation queries and directly consumable by BI tools such as Power BI. A star
schema is used in this project, implemented in SQLite for portability. The same SQL runs on PostgreSQL or MySQL.
\citet{kleppmann2017} discusses the wider design of data systems, in particular the separation of
transaction-processing stores from analytical stores and the use of simple, stable file formats between
processing stages. The framework follows this principle between its layers.

\section{Demand forecasting}
\subsection{Statistical methods}
Exponential smoothing dates back to the work of Holt and of \citet{winters1960}, who extended it to series with
trend and seasonality. The Holt--Winters method and its modern state-space formulation remain strong benchmarks.
So does the seasonal ARIMA (SARIMA) family, whose identification, estimation and diagnostic-checking cycle was
set out by \citet{box2015} \citep[see also][]{hyndman2021}. Both model a single time series and give
well-calibrated prediction intervals, but they cannot easily use additional explanatory variables or learn
shared patterns across many related series.

\subsection{Machine-learning methods}
Machine-learning regressors treat forecasting as supervised learning. Lagged values, rolling statistics and
calendar variables become features, and one global model can be trained across many series at once.
The simplest such learner is linear regression, whose estimation, diagnostics and assumptions are treated
comprehensively by \citet{montgomery2021}. Decision trees, formalised in the CART methodology of
\citet{breiman1984}, split the data recursively into homogeneous groups and form the building block of the
ensembles. Tree ensembles are the most common choice. Random forests average many decorrelated decision trees trained on
bootstrap samples \citep{breiman2001}. Gradient boosting fits trees sequentially to the residuals of the current
ensemble \citep{friedman2001}. XGBoost is a regularised and highly optimised implementation of gradient
boosting \citep{chen2016}. Support vector machines and support vector regression fit a maximum-margin function
in a kernel-induced feature space \citep{cortes1995, scholkopf2002}. $k$-nearest neighbours predicts from the
most similar historical cases, a method whose error was shown by \citet{cover1967} to be at most twice the
Bayes error rate as the sample grows. \citet{hastie2009} and \citet{james2021} provide the theory behind these
methods, particularly the bias--variance trade-off that motivates regularisation and hyper-parameter tuning.
\citet{bishop2006} and \citet{murphy2022} give the probabilistic view of the same models, while
\citet{kelleher2020}, \citet{geron2022} and \citet{raschka2022} show how they are applied in practice, with
worked predictive-analytics case studies and scikit-learn implementations. These books were the main guides for
the algorithm-level study in \cref{ch:algorithms}.

\subsection{Statistical versus machine-learning forecasting}
\citet{makridakis2018} compared eight statistical and ten machine-learning methods on 1,045 monthly series.
They found that, on average, simple statistical methods were more accurate than the machine-learning ones. They
concluded that machine-learning forecasts must always be benchmarked against simple methods. Two consequences
follow for this project. First, naive and seasonal-naive forecasts, together with Holt--Winters and SARIMA, are
included as benchmarks. Second, the evaluation design follows forecasting best practice. Time-series data must
not be randomly shuffled for validation, because doing so lets the model ``see the future''. Rolling-origin or
expanding-window schemes are preferred \citep{bergmeir2012, hyndman2021}. Forecasts made for different
hierarchical levels (product versus company) can be made coherent through reconciliation
\citep{athanasopoulos2009}.

\subsection{Accuracy measures}
\citet{hyndman2006} review forecast accuracy measures. Scale-dependent measures (MAE, RMSE) are interpretable
in the units of the data. Percentage errors (MAPE) are scale-free but explode when actual values are close to
zero. Measures relative to a benchmark show whether a method adds value. This study therefore reports MAE,
MSE, RMSE, $R^2$, MAPE and the weighted absolute percentage error (WAPE), together with the improvement over
naive baselines.

\section{Classification and evaluation}
Classification predicts a categorical outcome. Logistic regression remains a strong, interpretable baseline:
its coefficients are log odds ratios, and \citet{hosmer2013} describe how to build, assess and interpret such
models. Tree ensembles often capture interactions better \citep{kuhn2013}. When the classes are imbalanced, as with
large orders in this project, accuracy is misleading. A model that always predicts the majority class can look
accurate while being useless. Precision, recall, the F1-score, the area under the ROC curve (ROC-AUC) and the
area under the precision--recall curve are more informative. The decision threshold should be chosen on
validation data to match the business cost of errors \citep{kuhn2013}.

\section{Clustering and segmentation}
Market segmentation is a foundation of marketing strategy \citep{kotler2016}. Data-driven segmentation usually
relies on clustering. $K$-means, introduced by \citet{macqueen1967}, partitions observations into $k$ groups by
minimising within-cluster variance. \citet{kaufman1990} and \citet{tan2019} discuss its assumptions and
alternatives, including partitioning around medoids and hierarchical methods.
Agglomerative (Ward) clustering builds a hierarchy of merges. Gaussian mixture models assume each cluster is a
multivariate normal distribution \citep{han2012}. Because clustering has no ground truth, internal validity
indices are used. These include the silhouette coefficient \citep{rousseeuw1987}, the Davies--Bouldin index
\citep{davies1979} and the Calinski--Harabasz criterion.

\section{Anomaly detection}
Anomaly detection identifies observations that deviate from expected behaviour. Possible causes are data-entry
errors, fraud or genuine but exceptional events. Classical approaches compute robust distances from a fitted
expectation. The Isolation Forest \citep{liu2008} instead isolates points by random recursive partitioning:
anomalies need fewer splits to be isolated, which gives an efficient and distribution-free score. It is a
natural complement to rule-based checks \citep{aggarwal2017}.

\section{Explainable machine learning}
Managers rarely trust a model they cannot understand. Permutation importance measures how much a model's error
increases when a feature is randomly shuffled \citep{breiman2001}. SHAP values, grounded in cooperative game
theory, attribute each individual prediction additively to the input features \citep{lundberg2017}.
\citet{molnar2022} surveys these and other interpretation methods and distinguishes models that are
interpretable by design (linear models, short trees) from post-hoc explanations of complex models. Both kinds
are used here to show which factors drive demand forecasts and large orders.

\section{Dashboards and visual decision support}
\citet{few2013} defines a dashboard as a visual display of the most important information needed to achieve one
or more objectives, arranged on a single screen so that it can be monitored at a glance. He warns against
decorative charts, 3-D effects, gauges and colour used without meaning. He recommends instead clear
comparisons, sparse ink and consistent encodings. \citet{tufte2001} expresses the same idea as maximising the
\emph{data-ink ratio}, and \citet{knaflic2015} shows how to turn a chart into a business message by removing
clutter and directing attention to the point that matters. These principles guided the design of the dashboard in
\cref{ch:dashboard}. Modern open-source tools allow interactive dashboards to be built directly in the analysis
language and published free of charge. Examples are the Python libraries pandas \citep{mckinney2022},
scikit-learn \citep{pedregosa2011}, Plotly and Streamlit, whose use for data science is described by
\citet{vanderplas2016}.

\section{Research gap}
The literature offers mature techniques for each component: data cleaning, warehousing, forecasting,
classification, clustering, anomaly detection and dashboards. Studies and teaching materials, however, usually
treat them in isolation. Many applied forecasting studies also omit naive benchmarks or validate with random
splits, which makes their accuracy claims hard to trust \citep{makridakis2018, kaufman2012}. Few works show an
\emph{end-to-end}, reproducible framework in which (i) preprocessing, modelling and reporting are linked
automatically, (ii) evaluation is strictly time-ordered with benchmarks, and (iii) the outputs reach managers as
explained predictions and recommendations rather than raw model scores. This project addresses that gap for a
consumer-goods sales context in Bangladesh, a market that is under-represented in the published analytics
literature.

\begin{table}[htbp]\centering\small
\caption{Summary of key literature and its use in this project}\label{tab:litsummary}
\begin{tabularx}{\linewidth}{@{}p{3.6cm}X@{}}
\toprule
\textbf{Source} & \textbf{Contribution used in this project} \\ \midrule
\citet{shearer2000} & CRISP-DM process model underlying the methodology \\
\citet{chen2012}; \citet{sharda2018} & Descriptive--predictive--prescriptive framing of BI\&A \\
\citet{garcia2015}; \citet{han2012} & Cleaning, imputation, outlier and transformation techniques \\
\citet{kaufman2012} & Leakage-safe feature and validation design \\
\citet{kimball2013} & Star-schema data store \\
\citet{hyndman2021}; \citet{winters1960} & Holt--Winters, SARIMA, forecast reconciliation \\
\citet{makridakis2018}; \citet{bergmeir2012} & Benchmarks and time-ordered validation \\
\citet{breiman2001}; \citet{friedman2001}; \citet{chen2016} & Random forest, gradient boosting, XGBoost \\
\citet{rousseeuw1987}; \citet{davies1979} & Cluster validity indices \\
\citet{liu2008} & Isolation Forest anomaly detection \\
\citet{lundberg2017} & SHAP explanations \\
\citet{few2013} & Dashboard design principles \\ \bottomrule
\end{tabularx}
\end{table}

\section{Core textbooks}\label{sec:books}
Alongside the research papers, the project rests on a body of textbooks. \Cref{tab:books} lists the 34 books
consulted, grouped by subject, with the part of this work that each one informed. They range from the
foundations of statistical learning and forecasting to the practical engineering of data systems and the
communication of results to managers.

\begingroup\small
\setlength{\LTcapwidth}{\linewidth}
\begin{longtable}{@{}r >{\raggedright\arraybackslash}p{6.6cm} >{\raggedright\arraybackslash}p{5.9cm}@{}}
\caption{Textbooks consulted and their use in the project}\label{tab:books}\\
\toprule
\textbf{\#} & \textbf{Book} & \textbf{Used for} \\ \midrule
\endfirsthead
\toprule
\textbf{\#} & \textbf{Book} & \textbf{Used for} \\ \midrule
\endhead
\bottomrule
\endfoot
\multicolumn{3}{@{}l}{\emph{Business analytics and decision-making}} \\
1 & \citet{davenport2007}, \emph{Competing on Analytics} & business case for analytics (\cref{ch:intro}) \\
2 & \citet{provost2013}, \emph{Data Science for Business} & data-analytic thinking, problem framing \\
3 & \citet{sharda2018}, \emph{Business Intelligence, Analytics, and Data Science} & descriptive--predictive--prescriptive framing \\
4 & \citet{kotler2016}, \emph{Marketing Management} & market segmentation (\cref{ch:models}) \\
\multicolumn{3}{@{}l}{\emph{Data mining and data preparation}} \\
5 & \citet{han2012}, \emph{Data Mining: Concepts and Techniques} & knowledge-discovery process, outlier rules \\
6 & \citet{tan2019}, \emph{Introduction to Data Mining} & data quality, clustering \\
7 & \citet{witten2017}, \emph{Data Mining: Practical Machine Learning Tools} & preprocessing and evaluation practice \\
8 & \citet{garcia2015}, \emph{Data Preprocessing in Data Mining} & missing values, imputation hierarchy (\cref{ch:system}) \\
9 & \citet{kuhn2019}, \emph{Feature Engineering and Selection} & leakage-safe features, encodings \\
10 & \citet{aggarwal2017}, \emph{Outlier Analysis} & outlier and anomaly detection \\
\multicolumn{3}{@{}l}{\emph{Statistical and machine learning}} \\
11 & \citet{hastie2009}, \emph{The Elements of Statistical Learning} & bias--variance, ensembles, theory \\
12 & \citet{james2021}, \emph{An Introduction to Statistical Learning} & regression, classification, resampling \\
13 & \citet{bishop2006}, \emph{Pattern Recognition and Machine Learning} & probabilistic view of classifiers and kernels \\
14 & \citet{murphy2022}, \emph{Probabilistic Machine Learning} & modern treatment of all supervised models \\
15 & \citet{kuhn2013}, \emph{Applied Predictive Modeling} & tuning, imbalanced classes, thresholds \\
16 & \citet{kelleher2020}, \emph{Fundamentals of ML for Predictive Data Analytics} & worked predictive-analytics cases (\cref{ch:algorithms}) \\
17 & \citet{geron2022}, \emph{Hands-On Machine Learning} & scikit-learn practice, model tuning \\
18 & \citet{raschka2022}, \emph{Machine Learning with PyTorch and Scikit-Learn} & pipelines, scaling, evaluation \\
19 & \citet{montgomery2021}, \emph{Introduction to Linear Regression Analysis} & linear regression and diagnostics \\
20 & \citet{hosmer2013}, \emph{Applied Logistic Regression} & logistic regression and odds ratios \\
21 & \citet{breiman1984}, \emph{Classification and Regression Trees} & decision trees (CART) \\
22 & \citet{scholkopf2002}, \emph{Learning with Kernels} & support vector machines and kernels \\
23 & \citet{kaufman1990}, \emph{Finding Groups in Data} & cluster analysis and silhouette \\
24 & \citet{molnar2022}, \emph{Interpretable Machine Learning} & feature importance, SHAP, explanations \\
\multicolumn{3}{@{}l}{\emph{Forecasting and econometrics}} \\
25 & \citet{hyndman2021}, \emph{Forecasting: Principles and Practice} & Holt--Winters, accuracy measures, reconciliation \\
26 & \citet{box2015}, \emph{Time Series Analysis: Forecasting and Control} & SARIMA modelling \\
27 & \citet{wooldridge2020}, \emph{Introductory Econometrics} & fixed-effects price elasticity (\cref{ch:eda}) \\
\multicolumn{3}{@{}l}{\emph{Data engineering and tools}} \\
28 & \citet{kimball2013}, \emph{The Data Warehouse Toolkit} & star schema (\cref{ch:system}) \\
29 & \citet{kleppmann2017}, \emph{Designing Data-Intensive Applications} & layered data-system design \\
30 & \citet{mckinney2022}, \emph{Python for Data Analysis} & pandas data handling \\
31 & \citet{vanderplas2016}, \emph{Python Data Science Handbook} & NumPy, pandas, Matplotlib, scikit-learn \\
\multicolumn{3}{@{}l}{\emph{Visualisation and communication}} \\
32 & \citet{few2013}, \emph{Information Dashboard Design} & dashboard design (\cref{ch:dashboard}) \\
33 & \citet{tufte2001}, \emph{The Visual Display of Quantitative Information} & data-ink ratio, honest charts \\
34 & \citet{knaflic2015}, \emph{Storytelling with Data} & presenting insights to managers \\
\end{longtable}
\endgroup
